\documentclass[12pt,a4paper]{article}
\usepackage[UTF8]{ctex}        % 中文支持
\usepackage{geometry}
\geometry{left=2.54cm,right=2.54cm,top=3.18cm,bottom=3.18cm}
\usepackage{enumerate}        % 列表格式
\usepackage{titlesec}         % 标题格式
\usepackage{xcolor}
\usepackage{indentfirst}
\setlength{\parindent}{2em}
\usepackage{enumitem}
\usepackage{lipsum}           % 生成报文，测试用
\usepackage{adjustbox}
\usepackage{array}
\usepackage[table]{xcolor}

% ===================== 自定义功能区 =====================
\newcommand{\textred}[1]{\textcolor{red}{#1}}  %红色字体
\newcommand{\textgreen}[1]{\textcolor{green}{#1}}  %绿色字体
\newcommand{\textblue}[1]{\textcolor{blue}{#1}}  %绿色字体
% \newcommand{\wordnote}[2]{#1\,\raisebox{0.25ex}{\scriptsize #2}}
\newcommand{\upmean}[1]{\raisebox{0.25ex}{\scriptsize #1}}




\begin{document}
{\lishu \textcolor{gray}{在我们所降生其间的伟大谜团面前，你我从未停止像好奇的孩童一样站立凝视。}\par}



\section{4 types of question}
% ===================== 正文 =====================

{\large simon写作 Task2}

4 types of question:

1. Discussion

Some people think that it is more effective for students to study in groups, while others believe that it is better for them to study alone.
Discuss both views and give your own opinion.

{\small \textred{Plan: Topic - study in groups or alone.\\
\hspace*{4.4em} Answer - sometimes better alone, usually better in a group}}

People have different views about the effectiveness of group study as opposed to working alone.
While there are some benifits to studing independently, I beleve that group work is usually more productive.



2. Opinion

Some people believe that unpaid community service should be a compulsory part of high school programmes.
To what extent do you agree or disagree?

{\small \textred{Plan: Topic - community service for all teenagers\\
\hspace*{4.4em} Answer - 3 choices: agree, disagree, or balanced opinion (when possible)}}

It is sometimes argued that high school students should be made to do some work in their local communities. I completely agree that this kind of scheme would be a good idea.




3. Problem and solution

Many criminals reoffend after they have been punished. Why do some people continue to commit crimes after they have been punished, and what measures can be taken to tackle this problem?

{\small \textred{Plan:  Topic - criminals reoffend\\
\hspace*{4.4em} Answer - several reasons, a variety of measures (governments, communities)}}

It is true that punishments do not always deter criminals from committing more crimes. There are various reasons why offenders repeatedly break the law, but governments could certainly take steps to address this issue.




4. Two-part question

As most people spend a major part of their adult life at work, job satisfaction is an important element of individual well-being. What factors contribute to job satisfaction? How realistic is the expectation of job satisfaction for all workers?

{\small \textred{Plan:  Topic - job satisfaction\\
\hspace*{4.4em} Answer - several factors, unrealistic / impossible}}

Work plays a central role in our lives, and we would all like to feel fulfilled professionally. While a variety of factors may lead to job satisfaction, it would be unrealistic to expect everyone to be happy at work.


\newpage
\section{main body paragraphs}
\subsection{method 1: Firstly, Secondly, Finally}

Some people believe that unpaid community service should be a compulsory part of high school programmes.
To what extent do you agree or disagree?

3-minute plan:

- disagree for several reasons

- school timetable is full, no time for community service

- students' work in other subjects would be affected

- teenagers might not want to do it (reluctant, no motivation)

eg:

There are several reasons why I would argue against having compulsory community service for secondary school students.
Firstly, the school curriculum is already full with important academic subjects, such as maths, science and languages.
For example, I remember having an extremely busy timetable when I was at high school, and it would not have been possible to add to it.
Secondly, students' performance in other subjects would be affected if valuable study time were taken by charity work or neighbourhood improvement schemes.
Finally, I believe that teenage students would be reluctant to take part in any programme of obligatory work, and this could lead to poor motivation and even bad behaviour.

\subsection{method 2: Idea, Explain, Example}

Some people believe that unpaid community service should be a compulsory（旁边标注：voluntary） part of high school programmes.

To what extent do you agree or disagree?

3-minute plan:

- voluntary (not compulsory) community service is positive

- students more motivated if they can choose

- gain work experience, self confidence, skills

- good for CVs, career, university admissions, employers

eg:

On the other hand, the opportunity to do voluntary community service could be extremely positive for high school students.

By making these programmes optional, schools would ensure that only motivated students took part.

These young people would gain valuable experience in an adult working environment, which could help to build their self confidence and enhance their skills.

Having such experience and skills on their CVs could greatly improve school leavers’ career prospects.

For example, a period of voluntary work experience might impress a university admissions officer or a future employer.

\textgreen{link words} {\tiny link words do not help your 'vocabulary' score}

\textred{topic vocabulary} {\tiny examiners want to see 'topic vocabulary'}

这是合并了绿色和红色标注的完整文本：

\textgreen{On the other hand}, the \textred{opportunity} to do \textred{voluntary community service} could be extremely positive for high school students. 
By making \textgreen{these programmes} \textred{optional}, schools would \textred{ensure that only motivated students took part}. 
\textgreen{These young people} would \textred{gain valuable experience in an adult working environment}, which could help to build \textgreen{their} self confidence and \textred{enhance their skills}. 
Having \textgreen{such experience and skills} on their CVs could greatly improve \textred{school leavers' career prospects}. 
\textgreen{For example}, a period of voluntary work experience might \textred{impress a university admissions officer or a future employer}.




\section{conclusion}

4 types

\subsection{Discussion (+ opinion)}

1. 题目原文

In many cities the use of video cameras in public places is being increased in order to reduce crime, but some people believe that these measures restrict our individual freedom.
Do the benefits of increased security outweigh the drawbacks?

2. 开头段（Introduction）

It is true that video surveillance has become commonplace in many cities in recent years. While I understand that critics may see this as an invasion of privacy, I believe that there are more benefits than drawbacks.

3. 结尾段（Conclusion）

In conclusion, I would argue that the advantages of using video security systems in public places do outweigh the disadvantages.



\newpage
In many countries schools have severe problems with student behaviour.  
What do you think are the causes of this?  
What solutions can you suggest?  

\textblue{Plan your essay structure (4 paragraphs)  }

1. Introduction: student behaviour in schools variety of reasons, steps can be taken to tackle\\[-3pt]  
It is true that the behaviour of school pupils in some parts of the world has been getting worse in recent years. 
There are a variety of possible reasons for this, but steps can definitely be taken to tackle the problem.

2. \textblue{Causes of bad student behaviour}

In my opinion, three main factors are to blame for the way young people behave at school nowadays. Firstly, modern parents tend to be too lenient or permissive. 
Many children become accustomed to getting whatever they want, and they find it difficult to accept the demands of teachers or the limits imposed on them by school rules. 
Secondly, if teachers cannot control their students, there must be an issue with the quality of classroom management training or support within schools. 
Finally, children are influenced by the behaviour of celebrities, many of whom set the example that success can be achieved without finishing school.

3. \textblue{My suggested solutions}  

Student behaviour can certainly be improved. 
I believe that the change must start with parents, who need to be persuaded that it is important to set firm rules for their children. 
When children misbehave or break the rules, parents should use reasonable punishments to demonstrate that actions have consequences. 
Also, schools could play an important role in training both teachers and parents to use effective disciplinary techniques, and in improving the communication between both groups. 
At the same time, famous people, such as musicians and football players, need to understand the responsibility that they have to act as role models to children.

4. Conclusion: summarise the problem and steps

In conclusion, schools will continue to face discipline problems unless parents, teachers and public figures set clear rules and demonstrate the right behaviour themselves.

\textred{Band 7-9 vocabulary}

- steps can be taken to tackle(v 解决) the problem  

- three main factors are to blame  

- parents tend to be too lenient or permissive  

- children become accustomed\upmean{(vt 使习惯于)} to  

- limits imposed\upmean{(vi 欺骗；利用 vt 强加；以···欺骗；征税)} on them  

- quality of classroom management  

- celebrities, famous people, public figures  

- set an example  

- set firm rules, reasonable punishments  

- play an important role in  

- effective disciplinary techniques  

- responsibility to act as role models  

- face discipline\upmean{(vt 训练，训导；处罚，惩罚 n 训练，锻炼，训导)} problems





%============================
\newpage
{\Large
IELTS writing task 2 ~~ Lesson 9 ~~ 2-part question}

\textred{QUESTION:}
News editors decide what to broadcast on television and what to print in newspapers. 
What factors do you think influence these decisions? 
Do we become used to bad news, and would it be better if more good news was reported?

\textblue{Plan ideas for the two main paragraphs (6 minutes)}

1. Introduction: decisions about news stories
variety of factors, yes too much bad news

It is true that editors have to make difficult decisions about which news stories they broadcast or publish, 
and their choices are no doubt influenced by a variety of factors. 
In my opinion, we are exposed to too much bad news, and I would welcome a greater emphasis on good news.

2. Factors that influence news editors

Editors face a range of considerations when deciding what news stories to focus on. 
Firstly, I imagine that they have to consider whether viewers or readers will be interested enough to choose their television channel or their newspaper over competing providers. 
Secondly, news editors have a responsibility to inform the public about important events and issues, and they should therefore prioritise stories that are in the public interest. 
Finally, editors are probably under some pressure from the owners who employ them. For example, a newspaper owner might have particular political views that he or she wants to promote.

3. Too much bad news, should report more good

It seems to me that people do become accustomed to negative news. 
We are exposed on a daily basis to stories about war, crime, natural disasters and tragic human suffering around the world. 
I believe that such repeated exposure gradually desensitises people, and we become more cynical about the world. 
I would prefer to see more positive news stories, such as reports of the work of medical staff after a natural disaster, or the kindness of volunteers who help in their communities. 
This kind of news might inspire us all to lead better lives.

4. Conclusion: difficult news choices, more positive

In conclusion, it must be extremely difficult for editors to choose which news stories to present, 
but I would like to see a more positive approach to this vital public service.

\textred{Band 7-9 vocabulary}

- exposed to bad news

- welcome a greater emphasis \upmean{n.强调,重点} on good news

- Editors face a range of considerations

- prioritise \upmean{v 按重要性列出﹔确定〔事项﹑问题等〕的优先顺序:} stories that are in the public interest

- under some pressure from the owners

- promote particular political views

- become accustomed to negative news

- natural disasters and tragic \upmean{adj. 悲剧的；悲痛的，不幸的} human suffering

- exposure \upmean{n. 暴露,显露;揭发，揭露} gradually desensitises \upmean{vt. 使```无感觉或不敏感；使脱敏} people

- cynical \upmean{adj. 愤世嫉俗的；冷嘲的} and sceptical

- inspire \upmean{vt. 鼓舞；激发；启示；产生；使生灵感} us all to lead better lives

- a more positive approach

- vital public service

vital \upmean{adj. 极重要的，必不可少的
维持生命所必需的；生命的
有生气的，充满生机的}



\newpage
{\Large
IELTS writing task 2 ~~ Lesson 10 ~~ Extra help for "agree or disagree?" questions}

\textred{QUESTION:}
The money spent by governments on space programmes would be better spent on vital public services such as schools and hospitals.  
To what extent do you agree or disagree?

Governments in some countries spend large amounts of money on space exploration programmes. 
I completely agree with the idea that these are a waste of money, 
and that the funds should be allocated to public services.

2. explain why ‘space’ spending should be stopped

3. explain why public service spending is better  

4. Conclusion: spend on services that benefit us all

\rule{\linewidth}{0.4pt}

\textred{QUESTION:}
Some people believe that hobbies need to be difficult to be enjoyable.  
To what extent do you agree or disagree?  

Some hobbies are relatively easy, while others present more of a challenge. 
Personally, I believe that both types of hobby can be fun, 
and I therefore disagree with the statement that hobbies need to be difficult in order to be enjoyable.

2. explain why easy hobbies can be enjoyable  

3. explain why difficult hobbies can be fun  

4. Conclusion: disagree that difficult hobbies are better





\rule{\linewidth}{0.4pt}

\textred{QUESTION:}
Many people say that we now live in 'consumer societies' where money and possessions are given too much importance.

To what extent do you agree or disagree?

It is true that many people criticise modern society because it seems to be too materialistic. 
I agree with this to some extent, 
but I do not think it is the case that everyone is a victim of consumer culture.

2. I believe many people do focus too much on money

3. However, many others are not money oriented

4. Conclusion: partly agree




\rule{\linewidth}{0.4pt}

\textred{QUESTION:}
In the last century, the first man to walk on the moon said it was "a giant leap for mankind". 
However, some people think it has made little difference to our daily lives. 
To what extent do you agree or disagree?

It is often argued that the act of sending a man to the moon has been of no benefit to normal people. While I agree that this is true in practical terms, I believe that the psychological impact of this great achievement should not be underestimated.

2. no benefit in practical terms (standard of living, health)

3. but it was an inspiring achievement

4. Conclusion: partly agree





\newpage
\section*{IELTS Speaking}

\begin{itemize}
    \item 3 part of the test
    \item scoring system
    \item some overall advice 
\end{itemize}
\rule{\linewidth}{1pt}


Timing, 3 parts of the test

\begin{itemize}
    \item 11 to 14 minutes overall
    \item timing is very strict
    
    \item Part 1: 4 to 5 minutes
    \item Part 2: 3 to 4 minutes
    \item Part 3: 4 to 5 minutes
\end{itemize}
\rule{\linewidth}{1pt}


{\Large We need to work on 2 things: \par}

\begin{enumerate}
    \item \textred{Exam techniques / methods (use this video course)}
    \item \textred{Your gengeral level of English (use my blog, books, TV, films, dictionary, newspapers)}
\end{enumerate}

\rule{\linewidth}{1pt}
\textbf{Don't worry about:}
\begin{itemize}[leftmargin=*]
    \item body language
    \item eye contact
    \item the quality of your ideas
    \item whether the examiner agree with you
\end{itemize}

\rule{\linewidth}{1pt}
\textbf{Aims of this course:}
\begin{itemize}[leftmargin=*]
    \item give you a method
    \item prepare topics and common questions
\end{itemize}

\rule{\linewidth}{1pt}
Some tips:
\begin{itemize}
    \item natural language, not difficult language
    \item Answer as quickly as you can
\end{itemize}

\rule{\linewidth}{2pt}


A good way to practice:

\textit{ Do you like ...? }

\textblue{ games, walking, gardens, flowers}

\subsection*{IELTS Speaking part 2}

- 3 to 4 minute

- 1 minute to prepare

- speak for 2 minutes

- task card

Make notes:

Quickly decide what you're going to talk about

Technique:

Follow the points on the task card. Tery help you 
structure and give you something to say. Say as 
much as you can for each point.

Before the exam:

Prepare for common topic areas.

Vocabulary and ideas are the key.

\subsection*{Speaking lesson 3}

\textcolor{gray}{Topics, techniques, vocabulary, example qiestion}

Topic - 6 mian topic areas:

Describe a person/a place/ an object/ an event/ 
an activity/ your favourites. 

Prepare ideas and vocabulary (not grammar or linking)

Speaking naturally, explain in detail (content, not structure 结构没有很重要)

\subsubsection*{Describe a person}

涉及的职业：

teacher, famous person, firend, family member, 
child, someone who helps people, 
someone who does something well.

“他”的详细信息：

who (appearance, personality), \par
what he/ she does, \par
when/ how/ where you met, \par 
how you know about this person, \par
why you like this person. \par

\textred{Prepare tips:}

- think of a "\textred{theme} {\footnotesize hard-working}" for any person 
(firend, family, famous, child etc.)

- start with easy \textred{adjectives} {\footnotesize busy, active}

- search for \textred{better words} {\footnotesize conscientious, someone you can count on} and phrases

\vspace{0.54cm}
\textbf{hard-working ‘theme’:}

- enthusiastic  
- energetic  
- studious  
- persistent  
- motivated  
- determined to succeed  
- someone who sees things through  
- a good team player  
- likes to challenge himself / herself

根据图片内容，提取的文字如下：

\textbf{friendly ‘theme’:}

- kind, caring
- generous, unselfish
- big-hearted
- supportive
- down-to-earth, easy-going
- always there when you need him / her
- someone who cheers me up
- a big / magnetic personality
- lights up the room

Remember: 

You can't prepare for everything

\textred{you will need to adapt your ideas and improvise in the test.}

\subsubsection*{Describe a place}

Describe a place

- a city (you’ve visited)  

- a city (you would like to visit)  

- somewhere you went on holiday  

- a historic place  

- a river, lake or sea  

- a journey (where you went)  


- a shop  
- a restaurant  
- a street market



\newpage
\section*{Reading}
\subsection*{L1}

IELTS Reading: 1 hour, 40 questions.

Reading comprehension:Testing your understanding of written English, 
A vocabulary test.


\subsection*{L2}

{\centering \textcolor{blue}{Basic exam techniques:}\par}

\begin{enumerate}[label=\arabic*., topsep=0pt]
    \item Don't read the passage first
    \item Just read the title
    \item Then go to the first question
    \item Underline `keywords' in the question
    \item Then search for those words in the passage
    \item Underline them
    \item Read that part of the passage carefully
    \item Try to get the answer
\end{enumerate}

\vspace{1em}

\begin{enumerate}[label=\arabic*., topsep=0pt]
    \item Read at normal speed
    \item Only skim / scan for names or numbers
    \item Miss any difficult questions, \textcolor{red}{get to the end}
\end{enumerate}


\subsubsection*{Keyword table}

\rowcolors{2}{white}{white} % 行颜色交替（这里保持全白，方便你修改）
\begin{tabular}{|>{\raggedright\arraybackslash}p{5cm}|>{\raggedright\arraybackslash}p{8cm}|}
\hline
\rowcolor{purple!20} % 表头浅紫色背景
\textbf{Keywords in questions} & \textbf{Similar words in the passage} \\
\hline
inside & containing \\
\hline
shaded & shady \\
\hline
moist & damp \\
\hline
not classed as parasites & never parasitic \\
\hline
seen as & considered \\
\hline
special qualities & sense of calm, age and stillness \\
\hline
to dress wounds & as dressings on soldiers' wounds \\
\hline
\end{tabular}




























\newpage
\section*{近义词}
\begin{enumerate}
    \item spoil=lenient=permissive
    \item celebrity=famous people
    \item tackle=solve
\end{enumerate}






%============= 功能演示区 ============

% 单词\textsuperscript{释义}

% apple\textsuperscript{苹果} banana\textsuperscript{香蕉}


% study\upmean{学习} happy\upmean{开心}

\newpage


人类所有的不快乐，均来自一个原因：不知道如何安静地待在家里。

{\raggedleft —— 布莱斯·帕斯卡尔\par}













\end{document}